\section{Procedure}

The experimental apparatus and procedure followed the standard method
    outlined in the University of Victoria PHYS 326 laboratory manual~\cite{phys326}.

A circuit comprised of an inductor, resistor, and capacitor (LRC circuit) was assembled in series connected to a
function generator.
Measurements of the inductance, resistance, and capacitance were taken of each component with the digital multimeter
before connecting the circuit.
Table~\ref{tab:params} provides a list of measurements taken.
Uncertainty was calculated using the digital multimeter's spec sheet, including the range's relative and per digit
uncertainty.

\begin{table}[htbp]
\caption{Experimental Parameters}
\begin{tabular}{lrrr}
\toprule
{Quantity} & {Reading} \\
\midrule
$R_L$, Inductor resistance & \qty{109.4(1.3)}{\ohm} \\
L, Inductance & \qty{1.044(0.072)}{\henry} \\
C, Capacitance & \qty{10.07(0.21)}{\micro\farad}\\
$R_{FG}$, function generator resistance & \qty{50.00(0.25)}{\ohm} \\
$R_{res}$, normal & \qty{1133}{\ohm} \\
$R_{res}$, critical & \qty{485}{\ohm} \\
$R_{res}$, oscillatory & \qty{20}{\ohm} \\
\bottomrule
\end{tabular}\label{tab:params}
\end{table}

The function generator was configured to output a 5 Hz square wave signal to act as a continuous
switch for the LRC circuit to quickly cycle between charge and discharge.
The function generator offset was manually adjusted such that the ``ON'' and ``OFF'' settings were set to 3.86 V and
0 V, respectively.
Vernier differential probes were connected to the LabPro sensor interface and the LabPro software was calibrated to
take voltage readings at 3000 samples per second.

Before each data were collected, care was taken to ensure the probes were both physically shorted and zeroed in the
software to mitigate charge buildup that might offset the readings.
A preliminary test case was done to adjust the maximum recording time so that the majority of the discharge cycle
was captured without extending to the charging cycle.
The maximum time was set to 0.09 seconds, which appeared to consistently capture the whole event.
Resistor values were calculated to define damped, critically damped, and oscillatory behaviour for the LRC circuit.
The critical resistance value for the damped case was set to be double that of the critically damped case to ensure
signal behaviour between the two cases were appropriately different.
Once the decade resistance box was set using the dials, a discharge cycle was recorded via a voltage trigger and values
for inductor,
resistor, and
capacitor voltages produced.
This was done a total of three times, completed once for each damped, critically damped, and oscillatory case.

